\subsection{多変数関数の微分}

\subsubsection{偏微分}

2変数関数 $z=f(x,y)$ に対し、一方の変数を固定して他方だけを
変化させた微分を\textbf{偏微分}という。
\begin{align}
 f_x(x,y)&=\lim_{h\to0}\frac{f(x+h,y)-f(x,y)}h,\\
 f_y(x,y)&=\lim_{h\to0}\frac{f(x,y+h)-f(x,y)}h.
\end{align}
$f_x=\partial f/\partial x$ などと書く。高階偏導関数は
\[
 f_{xx}=\frac{\partial^2f}{\partial x^2},\quad
 f_{xy}=\frac{\partial^2f}{\partial y\partial x},\quad
 f_{yx}=\frac{\partial^2f}{\partial x\partial y}
\]
と表す。混合偏導関数が近傍で連続なら $f_{xy}=f_{yx}$ である。

\begin{center}
\begin{tikzpicture}[>=Stealth,x={(.85cm,-.35cm)},y={(.85cm,.35cm)},z={(0cm,1cm)}]
 \draw[->] (0,0,0)--(4,0,0) node[below] {$x$};
 \draw[->] (0,0,0)--(0,4,0) node[right] {$y$};
 \draw[->] (0,0,0)--(0,0,3.2) node[above] {$z$};
 \foreach \i in {0,...,4}
  {\draw[gray] plot[domain=0:4,samples=20] (\x,\i,{.08*(\x-2)^2+.08*(\i-2)^2+1});
   \draw[gray] plot[domain=0:4,samples=20] (\i,\x,{.08*(\i-2)^2+.08*(\x-2)^2+1});}
 \fill (2.8,2.3,1.06) circle (1.6pt);
 \draw[blue,thick] (1.5,2.3,.95)--(4,2.3,1.35) node[right] {$x$ 方向};
 \draw[red,thick] (2.8,1,.95)--(2.8,3.8,1.35) node[right] {$y$ 方向};
\end{tikzpicture}

\small 曲面の座標方向に沿った断面の傾きが偏導関数である
\end{center}

\begin{example}
$f(x,y)=2x^2+3xy-y^2+4y^3$ なら
\[
 f_x=4x+3y,\qquad f_y=3x-2y+12y^2
\]
である。
\end{example}

\paragraph{演習}
$f(x,y)=2e^{2x}\sin\sqrt{xy}$ の $f_x,f_y$ を求めよ。ただし $xy>0$
とする。

\paragraph{解答}
$s=\sqrt{xy}$ とおくと
\[
 s_x=\frac{y}{2\sqrt{xy}},\qquad
 s_y=\frac{x}{2\sqrt{xy}}.
\]
積の微分と連鎖律により
\begin{align*}
 f_x
 &=4e^{2x}\sin s+2e^{2x}\cos s\,s_x\\
 &=4e^{2x}\sin\sqrt{xy}
   +e^{2x}\frac{y}{\sqrt{xy}}\cos\sqrt{xy},\\
 f_y
 &=2e^{2x}\cos s\,s_y\\
 &=e^{2x}\frac{x}{\sqrt{xy}}\cos\sqrt{xy}.
\end{align*}

\subsubsection{全微分と連鎖律}

$f(x,t)$ で $x=x(t)$ も $t$ に依存するとき、偏微分
$\partial f/\partial t$ は $x$ を固定した変化率である。一方、
合成関数 $f(x(t),t)$ の微分は
\begin{equation}
 \boxed{\frac{df}{dt}
 =\frac{\partial f}{\partial t}
 +\frac{\partial f}{\partial x}\frac{dx}{dt}}
\end{equation}
となる。一般に $z=f(x_1,\ldots,x_n)$ の全微分は
\[
 dz=\sum_{i=1}^n\frac{\partial f}{\partial x_i}\,dx_i
\]
である。

\begin{example}
$f(x,t)=at^2+btx,\ x(t)=\sin t$ とする。連鎖律により
\[
 \frac{df}{dt}=(2at+bx)+(bt)\cos t
 =2at+b\sin t+bt\cos t.
\]
直接 $f(t)=at^2+bt\sin t$ を微分しても同じ結果を得る。
\end{example}

\subsubsection{勾配と方向微分}

スカラー場 $f(x_1,\ldots,x_n)$ に対し
\begin{equation}
 \boxed{\nabla f
 =\left(\frac{\partial f}{\partial x_1},\ldots,
 \frac{\partial f}{\partial x_n}\right)^{\mathsf T}}
\end{equation}
を $f$ の\textbf{勾配}という。単位ベクトル $\vect{u}$ 方向の
方向微分は
\[
 D_{\vect{u}}f=\nabla f\mathbin{\boldsymbol{\cdot}}\vect{u}
 =\lVert\nabla f\rVert\cos\theta
\]
である。したがって $\nabla f$ は $f$ が最も急に増加する方向を向き、
その大きさは最大変化率である。また等高線・等位面に直交する。

\begin{center}
\begin{tikzpicture}[>=Stealth,scale=1.1]
 \draw[->] (-3.2,0)--(3.5,0) node[right] {$x$};
 \draw[->] (0,-2.2)--(0,2.5) node[above] {$y$};
 \foreach \r in {.7,1.2,1.8}{\draw[gray] (0,0) ellipse ({1.45*\r} and \r);}
 \coordinate (P) at (1.45,1.04);
 \fill (P) circle (1.7pt) node[below right] {$P$};
 \draw[->,very thick,blue] (P)--($(P)+(1.1,1.6)$)
   node[above] {$\nabla f(P)$};
 \draw ($(P)+(-.8,.55)$)--($(P)+(.8,-.55)$);
\end{tikzpicture}

\small 勾配は等高線の法線方向を向く
\end{center}
